% =====================================================================
%  Conferenze, Zagarolo (On The Road), giorno 1 — Talk 6
%  Da aspirante designer a lead: cinque miti sulla carriera
%  Fonte: trascrizione automatica "Track 14". Talk in inglese di Manos,
%  game designer con circa dieci anni di esperienza (Supermassive Games,
%  The Chinese Room), con domande dal pubblico; modera Claudia.
%  Etichette con prefisso "c2-6:". Nessun pacchetto aggiuntivo.
% =====================================================================

\chapter[Cinque miti sulla carriera del designer]{Da aspirante designer a lead:
cinque miti sulla carriera nel game design}
\label{cap:c2-6-miti-carriera}

L'ultimo intervento del primo giorno è un talk in inglese sulla carriera del game
designer, dall'ingresso nell'industria fino al ruolo di \emph{lead}. Il relatore lo
costruisce intorno a una serie di \textbf{miti} che, a suo dire, accompagnano chi
lavora nell'industria nelle varie fasi del percorso, e li smonta uno per uno con
aneddoti presi dalla propria esperienza. Il pubblico interviene spesso, anche a
metà talk. Il capitolo è organizzato così:
\begin{enumerate}
  \item il relatore e la sua seconda carriera (sezione~\ref{sec:c2-6-relatore});
  \item mito 1: basta mettere un piede nella porta
        (sezione~\ref{sec:c2-6-mito1});
  \item mito 2: bisogna impressionare con le idee
        (sezione~\ref{sec:c2-6-mito2});
  \item mito 3: ci penserà l'arte (sezione~\ref{sec:c2-6-mito3});
  \item mito 4: l'esperienza dà certezze (sezione~\ref{sec:c2-6-mito4});
  \item mito 5: i grandi designer diventano lead (sezione~\ref{sec:c2-6-mito5});
  \item le domande del pubblico (sezione~\ref{sec:c2-6-domande}).
\end{enumerate}

% ---------------------------------------------------------------------
\section{Il relatore}
\label{sec:c2-6-relatore}

Il relatore è \textbf{Manos}\footnote{La trascrizione lo chiama ``Manus'' e, alla
fine, ``Mano''; il moderatore del capitolo~\ref{cap:c2-4-nfts} lo nomina come
``Manos'', forma che adottiamo. Il cognome non è pronunciato.}, greco, game
designer da circa \textbf{dieci anni}. Apre con una battuta: vuole fotografare la
sala e mandare la foto a \textbf{Rockstar}, per mostrare quante persone hanno
preferito la conferenza al nuovo \emph{GTA}, uscito proprio quel giorno.

\subsection{Una seconda carriera}

Prima dei videogiochi ha fatto ``molte altre cose'':
\begin{itemize}
  \item una prima laurea in \textbf{linguistica e letteratura};
  \item un lavoro all'\textbf{università};
  \item poi un impiego nella \textbf{pubblica amministrazione}, in
        un'organizzazione grande e nota.
\end{itemize}
Una sera, bevendo qualcosa con i colleghi dopo il lavoro, uno di loro gli disse che
le sue \textbf{uniche} aspirazioni nella vita erano \textbf{potere e denaro}. Manos
precisa di non avere nulla contro chi aspira a potere e denaro, aspirazioni
``universali''; ma il fatto che fossero le \emph{uniche} lo fece sentire fuori
posto. Aveva \textbf{32 anni} e si chiese: che cosa ci faccio qui? Voglio restare in
questo ambiente? La risposta fu no.

Sul treno di ritorno lesse per caso, su una rivista\footnote{La trascrizione
riporta ``Earth Magazine''; potrebbe trattarsi di \emph{Edge}, la nota rivista
britannica di videogiochi, ma è un'ipotesi.}, un articolo su un corso piuttosto
nuovo di una scuola di cinema britannica\footnote{Nella trascrizione ``an
abstract film school'': si tratta quasi certamente della \textbf{National Film and
Television School} (NFTS), presentata nel capitolo~\ref{cap:c2-4-nfts}, dove
Manos ha studiato.}. Pensò: ``qui posso fare giochi'', cosa che quando era cresciuto
e aveva studiato in \textbf{Grecia} non era considerata una carriera possibile. Si
iscrisse, tornò a scuola e si \textbf{indebitò di nuovo}, dopo aver appena finito
di ripagare il primo debito universitario. Da allora lavora nei videogiochi.

Il talk racconta il percorso ``suo, ma in fondo di tutti'' da \textbf{aspirante
designer} a \textbf{lead}, attraverso i miti che ha riconosciuto nella testa delle
persone nelle diverse fasi della carriera.

% ---------------------------------------------------------------------
\section{Mito 1: ``basta mettere un piede nella porta''}
\label{sec:c2-6-mito1}

Il primo mito è che per entrare nell'industria qualunque ruolo vada bene, purché
si entri. Manos dice di essere in parte in disaccordo con
Tony\footnote{Tony Evans, il producer del capitolo~\ref{cap:c2-2-producer}, che
aveva suggerito di muoversi in modo tattico anche accettando un primo ruolo
diverso da quello desiderato.}: è vero che il primo ruolo spesso non è esattamente
quello che si vuole, ma \textbf{per quanto possibile} bisognerebbe fare il più
possibile di \textbf{ciò che si vuole davvero fare}.

\subsection{La storia di Tom}

Alcuni anni prima il suo team cercava un \textbf{level designer}. Si presentò un
candidato, \textbf{Tom}, con un portfolio molto buono: un ottimo designer. Ma
durante il colloquio era evidente che la sua passione non era il level design ma il
\textbf{technical design}: si vedeva dal portfolio e dal modo in cui formulava le
risposte. Era preparato a rispondere sul level design, ma tutto faceva pensare che
non fosse ciò che voleva fare.

Manos, che ammette di avere la fama di ``non filtrarsi'', glielo chiese
apertamente: vuoi davvero fare il level designer? Tutto quello che mi hai detto
porta altrove; perché ti sei candidato? La risposta: il suo docente all'università
aveva detto alla classe che il level design era \textbf{una delle due principali
porte d'ingresso} nell'industria. ``Una sciocchezza'', commenta Manos: perché
qualcuno dovrebbe dirlo, e perché qualcuno dovrebbe crederci?

Gli disse che l'avrebbero considerato comunque, perché come level designer si era
presentato bene, ma gli consigliò di \textbf{ritirare la candidatura} se voleva
fare il technical designer: prima o poi sarebbe uscita una posizione adatta. Tom
accettò. \textbf{Cinque mesi dopo} si aprì una posizione da technical designer e lo
chiamarono direttamente: ``ora c'è un'apertura per quello che ci hai detto di voler
fare davvero''. È diventato \textbf{uno dei migliori junior} con cui Manos abbia
lavorato.

Questo non significa che nel lavoro si farà sempre ciò per cui si è stati assunti:
il lavoro reale è spesso molto diverso da quello immaginato. Ma
\textbf{idealmente}, quando si comincia un lavoro, bisognerebbe tenere davvero alle
cose che si faranno ogni giorno.

\subsection{Domanda: come capire il ruolo giusto?}

Dal pubblico: non è facile, da soli, capire per quale ruolo candidarsi e
riconoscere i propri punti di forza. Manos risponde che chi si candida deve
\textbf{fare i compiti}: capire di che cosa si occupa il ruolo, a partire dalla
descrizione dell'offerta.

Con l'IA, racconta, valutare le candidature è diventato ``un incubo'': domande
perfette sulla carta, poi al colloquio persone senza interesse o senza conoscenza
del ruolo. Sei o sette mesi prima avevano aperto una posizione da \textbf{combat
designer}: circa \textbf{150 candidature} nella prima settimana, in larga parte
dignitose (linguaggio giusto, lettere di presentazione decenti, buoni esempi). Ne
hanno chiamate a colloquio \textbf{otto}, ed è stato ``scioccante quanto fossero
scarsi'' alcuni: non solo mancavano le competenze, mancava \textbf{qualunque
interesse} per il combat design. Una candidata, a fronte di risposte di quattro o
cinque parole sulle domande di combat design, ha parlato per circa sette minuti di
fila di un argomento del tutto estraneo\footnote{La trascrizione riporta ``spoke
straight for about seven minutes about her dumb'': l'argomento non è
ricostruibile (forse ``her dog'', il suo cane).}.

Il consiglio: bisogna impegnarsi in entrambe le direzioni, scrivendo le candidature
in prima persona e informandosi su ciò per cui ci si candida. Fuori dall'industria
è difficile sapere com'è davvero il lavoro quotidiano, ma bisogna almeno provare a
farsene un'idea; e qui l'IA può essere utile: invece di farle scrivere la
candidatura, \textbf{chiederle che cosa fa quella disciplina}.

% ---------------------------------------------------------------------
\section{Mito 2: ``bisogna impressionare con le idee''}
\label{sec:c2-6-mito2}

Il secondo mito è che un designer debba \textbf{impressionare con le proprie
idee}. Per Manos ``\textbf{le idee costano poco}'', e l'affermazione ha tre
problemi:
\begin{enumerate}
  \item voler impressionare incoraggia una \textbf{cultura del rumore}, in cui si
        \textbf{impongono} le proprie idee agli altri;
  \item le idee non sono l'unica cosa in cui un designer deve essere bravo;
  \item tratta la creatività come una \textbf{costante}, mentre il bisogno di
        creatività cambia nelle diverse fasi del progetto.
\end{enumerate}

\subsection{Creatività e stress nel tempo}

Manos mostra un grafico ``un po' per scherzo, ma come molti scherzi sensato'': la
\textbf{creatività richiesta} varia nel corso dello sviluppo, e lo stesso vale per
lo \textbf{stress}. Combinando le due curve, verso la \textbf{fine del progetto}
la creatività richiesta è \textbf{bassa} e lo stress \textbf{alto}: chi in quella
fase vuole ancora impressionare con nuove idee crea problemi.

\subsection{Gli aneddoti di \emph{The Inpatient}}

Tutti, ammette, entrano nell'industria con quella mentalità, lui compreso. Il suo
primo lavoro a \textbf{Supermassive Games}\footnote{Nella trascrizione anche
``Supermaster''.} fu su \textbf{\emph{The Inpatient}}\footnote{Nella trascrizione
``the Impatient''. \emph{The Inpatient} (2018) è un horror per PlayStation VR di
Supermassive Games.}, un progetto ideale per lui: \textbf{VR}, su cui aveva
lavorato alla NFTS, e \textbf{horror}, su cui aveva letteralmente un master.
Il progetto veniva riavviato, e lui pensava di entrare e impressionare tutti con le
sue idee e con le sue critiche al design. Non considerava che, anche in un reboot,
c'erano persone che ci pensavano da moltissimo tempo e che esisteva una
\textbf{lore} di cui lui, da consumatore, non sapeva nulla, ma di cui bisognava
tenere conto.

Poi ci sono gli esempi estremi. Nello stesso progetto, verso la fine dello
sviluppo, una persona passò sei o sette ore del proprio tempo a progettare di sua
iniziativa una sequenza finale ``che faceva venire la nausea'' e la mandò ai
direttori del progetto e dello studio. Fu \textbf{l'unica persona del team a cui
non fu rinnovato il contratto}\footnote{La trascrizione aggiunge ``He was a 10'',
non interpretabile con sicurezza (forse un riferimento al tipo di contratto, per
esempio a termine).}.

Il problema, precisa, non era che avesse usato tempo di lavoro: \textbf{da lead,
bisogna avere a cuore anche il tempo personale delle persone}. Il fatto che lo
facesse nel proprio tempo libero non è una scusa: restava due o tre ore dopo il
lavoro per una settimana, e anche questo è un problema.

% ---------------------------------------------------------------------
\section{Mito 3: ``ci penserà l'arte''}
\label{sec:c2-6-mito3}

Il terzo mito è ``\textbf{art will fix it}'', dove ``arte'' sta per
\textbf{qualunque altra disciplina}. È l'idea antiquata, presente a tutti i
livelli, che le discipline siano \textbf{separate}: io faccio il mio design, sarà
ottimo, poi l'arte (o il suono, o la UI) lo porterà al livello in cui tutti ne
capiranno la grandezza.

\subsection{Coinvolgere le discipline fin dall'inizio}

Manos ha lavorato con un lead che diceva ai level designer: ``non preoccupatevi
della \textbf{contestualizzazione}, se ne occuperà l'arte; pensate solo che il
livello sia divertente''. In parte è vero, ma non si crea tutto per poi passarlo
agli artisti sperando che traducano e contestualizzino ogni cosa: le altre
discipline vanno \textbf{coinvolte durante tutto il processo}.

La disciplina che ha più da offrire e che viene coinvolta sempre \textbf{troppo
tardi} è il \textbf{suono}: nessuno chiede al sound team che cosa potrebbe fare per
semplificare la vita agli altri, eppure ha spesso ottime idee per risolvere
problemi.

\subsection{Blockout e prototipi non mentono}

L'altra faccia del mito: \textbf{i blockout e i prototipi non mentono}. Se il
blockout di un livello è pessimo, nessuna quantità di polish lo renderà migliore:
lo farà \emph{sembrare} migliore, ma non lo farà \emph{giocare} meglio. Arte e
design devono lavorare insieme; sembra ovvio, ma è facilissimo
\textbf{chiudersi in un silo}, soprattutto sotto stress creativo e con scadenze
strette.

\subsection{Intervento dal pubblico: ritmo e tempo}

Un partecipante trova il mito ``brutalmente vero'': spesso si sente dire che
servono più UX e UI, più effetti, per dare vita a una meccanica; ma una meccanica
che non funziona con cubi e semplici particelle non funzionerà nemmeno con VFX e
suono bellissimi. Ciò che conta, dice, è il ``ritmo'' giusto, il momento giusto.

Manos traduce: si tratta di \textbf{pacing} e \textbf{flow}, che si possono
sviluppare con \textbf{strumenti semplicissimi}. Nei motori moderni si possono usare
luci e altri strumenti a disposizione del designer prima di coinvolgere altre
persone.

Porta un esempio da \textbf{The Chinese Room}\footnote{Nella trascrizione ``the
Chinese room'': lo studio britannico The Chinese Room.}: un progetto che una grande
azienda avrebbe dovuto finanziare, ma di cui poi tagliò quasi tutti i fondi, quindi
non andò in porto. Era in parte un'avventura e in parte un \textbf{gioco di
carte}\footnote{La trascrizione riporta ``part of the Quantum Click Adventure and
part of the Cargill'': il genere preciso non è ricostruibile; la seconda parte
sembra ``card game'', come conferma la frase successiva.}. Prototiparono il gioco di
carte con \textbf{asset già pronti}, presi da un sito con migliaia di elementi,
\textbf{stampandolo su carta} con parole semplici, ed era divertente. Il
bilanciamento sarebbe venuto dopo, ma il cuore c'era, ``nonostante le mie terribili
capacità grafiche'': bastavano poche carte fatte da lui.

% ---------------------------------------------------------------------
\section{Mito 4: ``l'esperienza dà certezze''}
\label{sec:c2-6-mito4}

\subsection{Che cosa si intende per esperienza}

L'esperienza, nell'industria, è misurata in due modi: il \textbf{tempo} e i
\textbf{traguardi} raggiunti. Le offerte per ruoli senior chiedono ``X anni
nell'industria e almeno un titolo pubblicato''\footnote{La trascrizione riporta
``achieved at least one time'' e, poco dopo, ``the first gate they ever skipped'':
quasi certamente ``shipped at least one title'' e ``the first game they ever
shipped''.}, oppure di aver attraversato \textbf{l'intero ciclo di vita} di un
progetto.

Manos lo trova a volte ingiusto, perché nell'industria \textbf{molti progetti
vengono cancellati}. Conosce persone che lavorano da circa \textbf{nove anni} e
hanno pubblicato il loro primo gioco solo l'anno prima, dopo una serie di progetti
cancellati per un motivo o per l'altro: per anni sono state escluse da molte
posizioni senza alcuna colpa.

\subsection{Il riso per sushi}

Molti credono che con l'esperienza si diventi \textbf{sicuri} dei propri design, e
che il design diventi pianificabile. Per Manos è ``una sciocchezza'': l'esperienza è
solo esperienza, e \textbf{ogni progetto è unico}. La sua metafora:
\begin{quote}
per cuocere il riso da sushi servono 1,2 tazze d'acqua per tazza di riso crudo. Per
due tazze di riso non servono 2,4 tazze d'acqua: a seconda della durezza
dell'acqua possono bastarne 2,3 o 2,2. Finché non hai provato a cuocerne più di una
tazza, non sai quanta ne serve.
\end{quote}
Con i giochi è lo stesso. Aver lavorato su un gioco, anche dello stesso tipo, dà
\textbf{esperienza trasferibile}; ma con un team diverso, meccaniche un po' diverse,
dimensioni diverse, si sta facendo \textbf{un altro gioco}. All'inizio si possono
fare ipotesi, con un grado di certezza più alto di chi non ha mai fatto un gioco
simile, ma il primo design non sarà buono.
\begin{quote}
L'esperienza è esperienza. La qualità viene dall'iterazione.
\end{quote}

% ---------------------------------------------------------------------
\section{Mito 5: ``i grandi designer diventano lead''}
\label{sec:c2-6-mito5}

\subsection{Due percorsi: lead e principal}

L'ultimo mito è che i \textbf{grandi designer} facciano grandi cose e che il
percorso naturale sia: junior o mid, poi senior, poi \textbf{lead}. Non è
necessariamente così: si può diventare anche \textbf{principal}, e i due percorsi
richiedono competenze molto diverse.
\begin{description}
  \item[Lead] richiede \textbf{soft skill} molto più forti delle competenze di
        mestiere.
  \item[Principal] richiede di mantenere e approfondire soprattutto le
        \textbf{competenze di mestiere} (\emph{craft skill}).
\end{description}

\subsection{La scala che si ferma}

Il problema dell'industria: in tutte le aziende in cui Manos ha lavorato, e nella
maggior parte di quelle che conosce, \textbf{la scala del principal si ferma lì}.
Il lead invece può diventare project director, design director e così via. Torna
il tema di \textbf{potere e denaro}: le persone vogliono guadagnare di più ed
essere riconosciute, e la \textbf{semantica del ruolo} conta. Anche con venti livelli
di principal e stipendi sempre più alti, la preferenza andrebbe comunque al ruolo
di lead, perché nel lavoro successivo il lead apre altri gradini, mentre il
principal resta principal ``in perpetuo''.

Manos ammette di non avere una soluzione. In uno studio in cui ha
lavorato\footnote{La trascrizione riporta ``in M2H'': il nome dello studio non è
ricostruibile con sicurezza.} esistevano dei \textbf{design manager} che toglievano
ai lead una parte del lavoro amministrativo, ma non è sicuro che sia l'approccio
migliore. Forse si potrebbero \textbf{modularizzare} le descrizioni dei ruoli, per
far crescere insieme persone con competenze complementari; non ha visto una
soluzione che funzioni ovunque.

\subsection{L'industria non insegna le soft skill}

``Grande designer'', poi, è un'espressione vaga: di solito significa bravo nel
proprio mestiere. Ma l'industria \textbf{non insegna le soft skill} necessarie a un
lead: si arriva a senior avendone sviluppate pochissime, al massimo il
\textbf{mentoring}. Ci sono cose molto difficili da imparare dentro un'azienda:
\begin{itemize}
  \item \textbf{orientarsi nella struttura organizzativa}, un'abilità vera e propria;
  \item sapere \textbf{quando opporsi e quando no}, soprattutto in un ambiente
        \textbf{multiculturale}.
\end{itemize}
Bisogna trovare un modo per permettere la crescita professionale \textbf{senza
dare per scontato} che la strada verso il ruolo di lead sia quella più
gratificante.

% ---------------------------------------------------------------------
\section{Le domande del pubblico}
\label{sec:c2-6-domande}

\subsection{Quando la carriera si blocca}

Domanda: hai mai avuto momenti di ``sviluppo arrestato'', in cui la carriera non
andava avanti? Come li hai superati? Manos dice di essere stato fortunato: dieci
anni di crescita abbastanza costante. Ma ammette, ``ed è una cosa terribile da
dire'', che spesso nell'industria il modo per superare un blocco è
\textbf{cambiare lavoro}: dentro un'organizzazione i percorsi di crescita sembrano
più chiusi per chi c'è già che per chi arriva da fuori.

Il suo consiglio è parlarne con il proprio \textbf{manager} alla prima occasione e
sostenere le proprie ragioni. Almeno nel Regno Unito, nota, le persone sono molto
restie a parlare apertamente del proprio lavoro, e la trova una cultura
``terribile''. Sentirsi bloccati può anche significare che il manager non vede
crescita in un aspetto del ruolo che si aspetterebbe, e potrebbe non dirlo mai.
Parlarne è buon senso, ma dovrebbe essere \textbf{sempre il primo passo}.

\subsection{Quando l'iterazione diventa troppa}

Domanda: se la qualità viene dall'iterazione, arriva un punto in cui iterare
aggiunge pochissimo e anzi peggiora le cose; ti è mai successo? Manos, vincolato
dagli \textbf{NDA}, risponde in generale: in quei casi il problema non è chi itera,
ma chi deve \textbf{guidare il progetto} e prendere le decisioni creative.

Per lui un grande lead è \textbf{invisibile}: guardando il progetto non si riesce a
indicare che cosa abbia fatto, ma ha saputo guidare le persone verso il lavoro
migliore. Cita un ``autore famoso''\footnote{Non nominato; la massima è
generalmente attribuita a Neil Gaiman.}:
\begin{quote}
quando qualcuno ti dice che qualcosa non va, di solito ha ragione; quando ti dice
che cosa non va esattamente e come aggiustarlo, di solito ha torto.
\end{quote}
Il lead deve accorgersi che qualcosa non va e \textbf{guidare chi itera} verso la
decisione giusta.

\subsection{Le soft skill più importanti}

Domanda: quali sono le soft skill più importanti per un lead? Manos ne indica
diverse.
\begin{itemize}
  \item \textbf{Un lead non deve cercare di farsi degli amici.} Deve tenere
        profondamente al proprio team e aiutarlo a ottenere ciò che gli serve, ma
        non è lì per diventare amico di tutti. È una delle soft skill più difficili
        da sviluppare.
  \item \textbf{Sapere quando opporsi e quando dire ``si fa così''.} Non sono i
        lead a decidere quando finisce un progetto o che cosa diventerà: di solito,
        a meno di un budget illimitato, decidono \textbf{tempo e denaro}, e gli
        \textbf{stakeholder}. I ritardi capitano per una ragione, e bisogna capire
        quando insistere e quando accettare di conviverci.
  \item \textbf{Orientarsi nelle strutture organizzative}, che più sono
        gerarchiche più diventano politiche, senza diventare meschini e senza
        restare invischiati nelle \textbf{fazioni}.
\end{itemize}

\subsection{Candidarsi al ruolo giusto in un mercato saturo}

Un partecipante obietta: con un mercato del lavoro così \textbf{saturo},
soprattutto per i ruoli entry level, il consiglio di candidarsi solo ai ruoli
desiderati è difficile da seguire; lui si è candidato a moltissime posizioni,
spesso senza ricevere né risposta né feedback sul portfolio.

Manos gli dà ragione: ``non ho alcun diritto di contraddirti, è la verità''. Ma
aggiunge due considerazioni.
\begin{enumerate}
  \item L'industria dei videogiochi \textbf{non è enorme}. Se ci si candida a un
        lavoro che non si vuole e per cui non si è preparati, chi fa il colloquio
        potrebbe non dare una seconda possibilità quando ci si candida, meglio
        preparati, per il ruolo giusto. Non per meschinità: con così tante
        candidature per ogni posizione, non c'è spazio.
  \item Il \textbf{networking}, di cui si è parlato molto nella giornata.
\end{enumerate}

\paragraph{Come ha ottenuto il primo lavoro.} Durante la NFTS Manos si era
candidato a uno stage a Supermassive, senza mai ricevere risposta, tre o quattro
mesi prima del diploma. Poi portò il suo gioco ad alcuni eventi. A uno di questi un
\textbf{writer a contratto} di Supermassive lo provò e gli disse che la sua azienda
avrebbe adorato quel tipo di design. Non aveva tempo di giocarlo per intero, ma si
sarebbero rivisti a un altro evento la settimana successiva: lì fece una partita
completa, chiese un'email con il link a una build e la passò all'azienda. Una
settimana dopo Manos ebbe un colloquio.

\paragraph{La fortuna.} Manos ammette che è stata \textbf{fortuna}, e che la
fortuna conta enormemente. Racconta un'intervista a \textbf{Christoph Waltz}, due
volte premio Oscar: alla domanda su quanto la fortuna avesse contato nella sua
carriera, si era sentito offeso (``ho vinto due Oscar, sono bravissimo in quello che
faccio''); poi, riflettendoci, aveva risposto che il \textbf{100\%} del suo
successo era dovuto alla fortuna, perché ci sono forse un milione di attori bravi
quanto lui che non si sono trovati nel posto giusto al momento giusto. Bisogna
\textbf{costruirsi la fortuna} il più possibile, mostrando al mondo ciò che si fa,
sapendo che qualcuno che lo fa meno potrebbe comunque avere successo prima.

\paragraph{Consigli pratici.} Il partecipante vorrebbe lavorare nel
\textbf{narrative design}, con anche un lato artistico, e non sa dove trovare
\textbf{portfolio review} in Italia. Manos suggerisce:
\begin{itemize}
  \item frequentare i luoghi dove si può mostrare il proprio lavoro;
  \item partecipare alle \textbf{game jam}, dove spesso ci sono persone
        dell'industria che possono vedere l'etica del lavoro e ciò che si produce
        sotto stress e in poco tempo;
  \item esplorare \textbf{combinazioni di competenze}: per esempio unire
        narrative design e level design (o l'arte) e presentarle come un'unica
        competenza vendibile.
\end{itemize}
Riconosce di parlare da una posizione privilegiata\footnote{La trascrizione
riporta ``sitting in my eyeball in Paris'', probabilmente ``sitting in my ivory
tower'': ``seduto nella mia torre d'avorio''.} e che il momento è difficile: ``sono
d'accordo con te''.

Il talk si chiude con i ringraziamenti\footnote{La trascrizione termina con
alcune decine di ripetizioni meccaniche di ``thank you'', omesse.}.

% ---------------------------------------------------------------------
\section*{Punti chiave del capitolo}
\addcontentsline{toc}{section}{Punti chiave del capitolo}

\begin{itemize}
  \item Candidarsi ai ruoli che si vogliono davvero fare e \textbf{informarsi} su
        che cosa comportano: un colloquio sbagliato può chiudere una porta in
        un'industria piccola.
  \item \textbf{Le idee costano poco}: la creatività richiesta cala verso la fine
        del progetto mentre lo stress cresce; bisogna rispettare la lore e il
        lavoro di chi c'era prima, e anche il proprio tempo libero.
  \item ``Ci penserà l'arte'' è un mito: coinvolgere tutte le discipline (il
        \textbf{suono} per primo) fin dall'inizio; \textbf{blockout e prototipi non
        mentono}, pacing e flow si costruiscono con strumenti semplici.
  \item L'esperienza non dà certezze (il riso per sushi): ogni gioco è diverso,
        \textbf{la qualità viene dall'iterazione}.
  \item Lead e principal sono percorsi diversi (soft skill contro craft skill);
        l'industria premia il ruolo di lead e non insegna le soft skill.
  \item Un buon lead è \textbf{invisibile}, non cerca amici, sa quando opporsi e
        sa orientarsi nell'organizzazione.
  \item Networking, eventi e game jam contano; la fortuna conta moltissimo, ma
        bisogna costruirsela mostrando il proprio lavoro.
\end{itemize}
